\section{A Framework for Self-Improving Intelligent Textbooks}

\subsection{The Core Analogy}

The mapping between DGM components and intelligent textbook components is
direct enough to be productive. Table~\ref{tab:analogy} makes it explicit.

\begin{table}[h]
\centering
\caption{DGM--Intelligent Textbook Component Mapping}
\begin{tabular}{@{}L{4.5cm}L{7.5cm}@{}}
\toprule
DGM Concept & Intelligent Textbook Analog \\
\midrule
Coding agent codebase & MicroSim source files and content-generation skills \\
Frozen foundation model & The code-generation language model (unchanged) \\
Benchmark (SWE-bench) & Student learning outcomes and engagement metrics \\
Archive of agents & Versioned library of MicroSim and explanation variants \\
Parent selection & Choosing a high-performing variant to mutate \\
Self-modification & Generating a content variant with a pedagogical hypothesis \\
Fitness evaluation & Aggregating analytics into a per-variant score \\
Open-ended exploration & Concept-graph traversal and Bloom's-level laddering \\
\bottomrule
\end{tabular}
\label{tab:analogy}
\end{table}

\subsection{The Evolutionary Loop}

The proposed system operates in a generational cycle analogous to the DGM's,
but with pedagogical artifacts as the unit of evolution:

\begin{enumerate}
    \item \textbf{Selection.} From the archive of content variants for a given
          concept node, select a parent. Selection favors variants with strong
          measured outcomes but reserves probability mass for under-explored
          variants, balancing exploitation against exploration just as the DGM
          balances high-performing parents against those with few children.

    \item \textbf{Mutation.} A code-generation tool produces one or more child
          variants of the parent. Each mutation is driven by an explicit
          \textit{pedagogical hypothesis}---for example, ``adding a real-time
          graph of velocity over time will improve student ability to connect
          the simulation to the underlying equation,'' or ``reducing the number
          of simultaneous parameters from four to two will reduce cognitive
          load for novices.''

    \item \textbf{Deployment.} Child variants are deployed to a subset of
          learners, either through randomized assignment (A/B testing) or
          through a multi-armed bandit allocation that shifts traffic toward
          better-performing variants over time.

    \item \textbf{Evaluation.} Learning analytics are gathered for each variant
          and aggregated into a fitness score (Section~\ref{sec:fitness}).

    \item \textbf{Archival.} Variants meeting a quality threshold are retained
          in the archive with their full performance record, becoming candidate
          parents for future generations. Crucially, even underperforming
          variants are retained as potential stepping stones, preserving the
          open-endedness that the DGM ablation studies showed to be essential.
\end{enumerate}

\subsection{The Concept Dependency Graph as Search Space}

In the DGM, the search space is the space of all computable agent programs.
In an intelligent textbook, the concept dependency graph provides a far more
structured and tractable search space. Each node is an independent
evolutionary population. The graph topology supplies three useful signals:

\begin{itemize}
    \item \textbf{Prioritization.} Concepts where students show weak mastery,
          or whose downstream dependents show degraded performance, are
          prioritized for evolutionary effort.
    \item \textbf{Diagnosis.} When learners consistently fail at a particular
          prerequisite edge, the system can hypothesize that a
          \textit{bridging concept} is missing and propose inserting a new
          node---a structural mutation of the graph itself, not merely of
          content within a node.
    \item \textbf{Credit assignment.} Because the graph encodes prerequisites,
          the system can partially attribute a student's failure on a downstream
          concept to weaknesses in upstream content, focusing improvement where
          it will have the greatest cascading benefit.
\end{itemize}

\subsection{The Fitness Function: From Analytics to Selection Pressure}
\label{sec:fitness}

The single most important design decision, as the AlphaEvolve experience makes
clear, is the evaluator. A naive fitness function---for example, raw quiz
score---invites perverse outcomes such as variants that make quizzes easier
rather than instruction better. A robust pedagogical fitness function should be
multi-dimensional and aligned with genuine learning.

We propose a composite fitness function with the following components:

\begin{itemize}
    \item \textbf{Learning gain.} The improvement in assessment performance
          from before to after engaging with the MicroSim, normalized by the
          maximum possible gain. This is the primary signal and should dominate.
    \item \textbf{Bloom's-level advancement.} Whether the student demonstrates
          mastery at higher cognitive levels (\textit{Apply}, \textit{Analyze})
          rather than mere recall. A variant that moves students up the taxonomy
          is worth more than one that merely improves recall.
    \item \textbf{Efficiency.} The number of attempts and the time required to
          reach mastery. Faster mastery, all else equal, indicates clearer
          instruction.
    \item \textbf{Engagement quality.} Interaction depth with the MicroSim,
          distinguishing meaningful exploration from disengagement or random
          clicking.
    \item \textbf{Retention.} Performance on the concept in later sessions,
          capturing durable learning rather than transient performance.
\end{itemize}

These components are combined into a single scalar with weights that reflect
institutional priorities, with learning gain and retention weighted most
heavily to guard against optimizing for shallow proxies. The composite design
directly echoes AlphaEvolve's requirement for an evaluation function with
explicit metrics, adapted to the multi-objective reality of learning.

\subsection{Two Deployment Phases}

A practical deployment must account for the reality that most intelligent
textbooks begin life on a static host such as GitHub Pages with anonymous
users, and only later integrate identity-bearing learning analytics. The
framework supports both phases.

\textbf{Phase~1: Anonymous aggregate analytics.} With static hosting and no
learner identity, the system relies on aggregate event analytics (for example,
a web analytics platform) and browser-local storage for within-session
continuity. This yields aggregate per-variant fitness signals sufficient to
drive selection at the population level, though it cannot reconstruct
individual learning trajectories.

\textbf{Phase~2: Identity-bearing xAPI and LRS.} As the textbook integrates
the Experience API and a Learning Record Store~\cite{adl2013xapi}, each
interaction is recorded as a structured statement carrying a learner
identifier, enabling per-learner trajectory analysis, retention measurement
across sessions, and far richer fitness signals.

The critical engineering recommendation is to \textbf{design the Phase~1
analytics schema to match the Phase~2 xAPI schema from the outset}, using
identical field names and concept identifiers. This makes the eventual
migration a matter of swapping the data sink rather than re-instrumenting
content, and allows fitness queries to run against either backend with minimal
change.
